\documentclass[11pt]{article}
\newcommand{\ListaTitulo}{Gabarito 12: Distribuições Uniforme e Normal}
% Preâmbulo compartilhado: MATF14, Listas de Exercícios 2026
% Usado por ListaNN.tex e GabaritoNN.tex via % Preâmbulo compartilhado: MATF14, Listas de Exercícios 2026
% Usado por ListaNN.tex e GabaritoNN.tex via % Preâmbulo compartilhado: MATF14, Listas de Exercícios 2026
% Usado por ListaNN.tex e GabaritoNN.tex via \input{preamble}

\usepackage[utf8]{inputenc}
\usepackage[T1]{fontenc}
\usepackage[brazil]{babel}
\usepackage{fullpage}
\usepackage{amsmath,amssymb}
\usepackage{enumitem}
\usepackage{xcolor}
\usepackage{fancyhdr}
\usepackage{titlesec}
\usepackage{listings}
\usepackage{hyperref}
\usepackage{booktabs}
\usepackage{graphicx}
\usepackage{tikz}

\definecolor{matf14azul}{HTML}{0D3B54}
\definecolor{matf14dourado}{HTML}{C9971E}

\titleformat{\section}{\normalfont\Large\bfseries\color{matf14azul}}{\thesection}{1em}{}
\titleformat{\subsection}{\normalfont\large\bfseries\color{matf14azul}}{\thesubsection}{1em}{}

\providecommand{\ListaTitulo}{Lista de Exercícios}

\setlength{\headheight}{14pt}
\pagestyle{fancy}
\fancyhf{}
\lhead{\small MATF14: Estatística Econômica I}
\rhead{\small \ListaTitulo}
\cfoot{\thepage}
\renewcommand{\headrulewidth}{0.4pt}

\lstset{
  basicstyle=\ttfamily\small,
  breaklines=true,
  frame=single,
  columns=fullflexible,
  backgroundcolor=\color{gray!8},
  commentstyle=\color{matf14dourado},
  keywordstyle=\color{matf14azul}\bfseries,
  language=R
}

\newcounter{questaocounter}
\newenvironment{questao}[1]{%
  \stepcounter{questaocounter}%
  \bigskip\noindent\textbf{Questão #1.}\ %
}{\par}

\newcommand{\cabecalho}[2]{%
  \begin{center}
    {\Large\bfseries\color{matf14azul} #1}\\[0.3em]
    {\large #2}
  \end{center}
  \vspace{0.5em}
  \noindent\rule{\linewidth}{0.6pt}
  \vspace{0.5em}
}

\newenvironment{solucao}{%
  \par\smallskip\noindent\textit{\color{matf14dourado!80!black}Solução.}\ %
}{\hfill$\blacksquare$\par}


\usepackage[utf8]{inputenc}
\usepackage[T1]{fontenc}
\usepackage[brazil]{babel}
\usepackage{fullpage}
\usepackage{amsmath,amssymb}
\usepackage{enumitem}
\usepackage{xcolor}
\usepackage{fancyhdr}
\usepackage{titlesec}
\usepackage{listings}
\usepackage{hyperref}
\usepackage{booktabs}
\usepackage{graphicx}
\usepackage{tikz}

\definecolor{matf14azul}{HTML}{0D3B54}
\definecolor{matf14dourado}{HTML}{C9971E}

\titleformat{\section}{\normalfont\Large\bfseries\color{matf14azul}}{\thesection}{1em}{}
\titleformat{\subsection}{\normalfont\large\bfseries\color{matf14azul}}{\thesubsection}{1em}{}

\providecommand{\ListaTitulo}{Lista de Exercícios}

\setlength{\headheight}{14pt}
\pagestyle{fancy}
\fancyhf{}
\lhead{\small MATF14: Estatística Econômica I}
\rhead{\small \ListaTitulo}
\cfoot{\thepage}
\renewcommand{\headrulewidth}{0.4pt}

\lstset{
  basicstyle=\ttfamily\small,
  breaklines=true,
  frame=single,
  columns=fullflexible,
  backgroundcolor=\color{gray!8},
  commentstyle=\color{matf14dourado},
  keywordstyle=\color{matf14azul}\bfseries,
  language=R
}

\newcounter{questaocounter}
\newenvironment{questao}[1]{%
  \stepcounter{questaocounter}%
  \bigskip\noindent\textbf{Questão #1.}\ %
}{\par}

\newcommand{\cabecalho}[2]{%
  \begin{center}
    {\Large\bfseries\color{matf14azul} #1}\\[0.3em]
    {\large #2}
  \end{center}
  \vspace{0.5em}
  \noindent\rule{\linewidth}{0.6pt}
  \vspace{0.5em}
}

\newenvironment{solucao}{%
  \par\smallskip\noindent\textit{\color{matf14dourado!80!black}Solução.}\ %
}{\hfill$\blacksquare$\par}


\usepackage[utf8]{inputenc}
\usepackage[T1]{fontenc}
\usepackage[brazil]{babel}
\usepackage{fullpage}
\usepackage{amsmath,amssymb}
\usepackage{enumitem}
\usepackage{xcolor}
\usepackage{fancyhdr}
\usepackage{titlesec}
\usepackage{listings}
\usepackage{hyperref}
\usepackage{booktabs}
\usepackage{graphicx}
\usepackage{tikz}

\definecolor{matf14azul}{HTML}{0D3B54}
\definecolor{matf14dourado}{HTML}{C9971E}

\titleformat{\section}{\normalfont\Large\bfseries\color{matf14azul}}{\thesection}{1em}{}
\titleformat{\subsection}{\normalfont\large\bfseries\color{matf14azul}}{\thesubsection}{1em}{}

\providecommand{\ListaTitulo}{Lista de Exercícios}

\setlength{\headheight}{14pt}
\pagestyle{fancy}
\fancyhf{}
\lhead{\small MATF14: Estatística Econômica I}
\rhead{\small \ListaTitulo}
\cfoot{\thepage}
\renewcommand{\headrulewidth}{0.4pt}

\lstset{
  basicstyle=\ttfamily\small,
  breaklines=true,
  frame=single,
  columns=fullflexible,
  backgroundcolor=\color{gray!8},
  commentstyle=\color{matf14dourado},
  keywordstyle=\color{matf14azul}\bfseries,
  language=R
}

\newcounter{questaocounter}
\newenvironment{questao}[1]{%
  \stepcounter{questaocounter}%
  \bigskip\noindent\textbf{Questão #1.}\ %
}{\par}

\newcommand{\cabecalho}[2]{%
  \begin{center}
    {\Large\bfseries\color{matf14azul} #1}\\[0.3em]
    {\large #2}
  \end{center}
  \vspace{0.5em}
  \noindent\rule{\linewidth}{0.6pt}
  \vspace{0.5em}
}

\newenvironment{solucao}{%
  \par\smallskip\noindent\textit{\color{matf14dourado!80!black}Solução.}\ %
}{\hfill$\blacksquare$\par}


\begin{document}

\cabecalho{Gabarito 12}{Distribuições Uniforme e Normal (Cap. 4, Aulas 27--28)}

\begin{questao}{1} \textbf{(Uniforme: horário de chegada de um ônibus)}
\begin{solucao}
\begin{enumerate}[label=(\alph*)]
  \item $E(X)=\frac{0+15}{2}=7{,}5$ minutos. Faz sentido: como qualquer instante entre dois
    ônibus é igualmente provável para a chegada do passageiro, em média ele chega "no meio" do
    intervalo, nem sempre logo após um ônibus (espera quase 15 min) nem sempre logo antes do
    próximo (espera quase 0).
  \item $P(X>10) = 1-F(10) = 1-\frac{10-0}{15-0} = 1-\frac{10}{15} = \frac13\approx0{,}333$.
  \item $dp(X) = \sqrt{\frac{(15-0)^2}{12}} = \sqrt{\frac{225}{12}} = \sqrt{18{,}75}\approx4{,}33$.
\end{enumerate}
\end{solucao}
\end{questao}

\begin{questao}{2} \textbf{(Normal: notas de uma prova)}
\begin{solucao}
\begin{enumerate}[label=(\alph*)]
  \item $Z=\frac{5{,}0-6{,}5}{1{,}2}=-1{,}25$. $P(X<5{,}0)=P(Z<-1{,}25)\approx0{,}1056$
    (10,56\%).
  \item Procura-se $x_0$ tal que $P(X>x_0)=0{,}10$, ou seja, $F(x_0)=0{,}90$. O percentil 90 da
    Normal padrão é $z\approx1{,}2816$. $x_0 = 6{,}5+1{,}2816\times1{,}2 \approx 6{,}5+1{,}538 =
    8{,}04$.
  \item Como a Normal é simétrica em torno da média, exatamente metade da distribuição fica
    abaixo da média: $P(X<6{,}5)=0{,}5$, logo esperam-se $200\times0{,}5=100$ alunos abaixo da
    média. Isso não deveria surpreender: é uma consequência direta da simetria, não uma
    coincidência dos dados, vale sempre para qualquer Normal, independente de $\mu$ e $\sigma$.
\end{enumerate}
\end{solucao}
\end{questao}

\begin{questao}{3} \textbf{(Padronização: comparando desempenho em escalas diferentes)}
\begin{solucao}
\begin{enumerate}[label=(\alph*)]
  \item Matemática: $z_{mat} = \frac{720-650}{80} = \frac{70}{80}=0{,}875$.
    Redação: $z_{red} = \frac{78-70}{10} = \frac{8}{10}=0{,}8$.
  \item O desempenho relativamente melhor foi em \textbf{Matemática} ($z=0{,}875 > 0{,}8$): o
    candidato ficou mais desvios padrão acima da média dos demais candidatos nessa prova do que
    na de redação, mesmo com notas brutas (720 vs. 78) em escalas completamente diferentes.
  \item As notas brutas estão em escalas com médias e desvios padrão diferentes entre as duas
    provas (650 pontos vs. 70 pontos de médias, 80 vs. 10 de desvios), comparar 720 e 78
    diretamente não diz nada sobre desempenho relativo; só a padronização (posição em unidades de
    desvio padrão frente à respectiva distribuição) permite uma comparação válida.
\end{enumerate}
\end{solucao}
\end{questao}

\begin{questao}{4} \textbf{(Regra empírica: controle de qualidade de peso)}
\begin{solucao}
\begin{enumerate}[label=(\alph*)]
  \item Pela regra empírica, $\approx95\%$ dos dados ficam entre $\mu\pm2\sigma =
    500\pm2(5) = [490, 510]$ gramas.
  \item $Z_1=\frac{490-500}{5}=-2$, $Z_2=\frac{510-500}{5}=2$.
    $P(490<X<510)=P(-2<Z<2)\approx0{,}9545$ (95,45\%), muito próximo dos 95\% aproximados da
    regra empírica (a regra é, de fato, uma aproximação arredondada desse valor exato).
  \item $Z=\frac{485-500}{5}=-3$. $P(X<485)=P(Z<-3)\approx0{,}00135$ (cerca de 0,135\%), um
    evento raro sob o modelo. Se a reclamação envolve pacotes "frequentemente" muito abaixo do
    rótulo, esse resultado \textbf{não sustenta} a reclamação como algo esperado do processo
    normal, sugere, em vez disso, que ou o processo real se afastou do modelo assumido
    (\emph{média} ou \emph{desvio padrão} mudaram), ou os casos reclamados são mesmo exceções
    raras, e não a regra.
\end{enumerate}
\end{solucao}
\end{questao}

\begin{questao}{5} \textbf{(Dedução: esperança e variância da Uniforme)}
\begin{solucao}
\begin{enumerate}[label=(\alph*)]
  \item
  $$
  E(X) = \int_a^b x\cdot\frac{1}{b-a}\,dx = \frac{1}{b-a}\left[\frac{x^2}{2}\right]_a^b
  = \frac{1}{b-a}\cdot\frac{b^2-a^2}{2} = \frac{(b-a)(b+a)}{2(b-a)} = \frac{a+b}{2},
  $$
  usando a fatoração $b^2-a^2=(b-a)(b+a)$.
  \item Primeiro, $E(X^2) = \int_a^b x^2\cdot\frac{1}{b-a}dx = \frac{1}{b-a}
  \left[\frac{x^3}{3}\right]_a^b = \frac{b^3-a^3}{3(b-a)}$. Usando a fatoração
  $b^3-a^3=(b-a)(b^2+ab+a^2)$: $E(X^2) = \frac{b^2+ab+a^2}{3}$. Logo,
  $$
  \mathrm{Var}(X) = E(X^2)-[E(X)]^2 = \frac{a^2+ab+b^2}{3} - \left(\frac{a+b}{2}\right)^2
  = \frac{a^2+ab+b^2}{3} - \frac{a^2+2ab+b^2}{4}.
  $$
  Colocando em denominador comum 12:
  $$
  \mathrm{Var}(X) = \frac{4(a^2+ab+b^2) - 3(a^2+2ab+b^2)}{12}
  = \frac{4a^2+4ab+4b^2-3a^2-6ab-3b^2}{12} = \frac{a^2-2ab+b^2}{12} = \frac{(b-a)^2}{12},
  $$
  como queríamos demonstrar.
\end{enumerate}
\end{solucao}
\end{questao}

\end{document}
